In particular, we can use this theorem to get the following result.

\begin{Theorem}
\label{thm:hilbert-schmidt-kernel-compactness}
Let the kernel $k:\Omega\times\Omega\to\mathbb{R}$ be Hilbert-Schmidt. Then the integral operator $K\in\mathcal{L}(L^2(\Omega))$ defined by\[K(u)(\bm{x}):=\int_\Omega k(\bm{x},\bm{y})u(\bm{y})\,\mathrm{d}\bm{y}\]is compact.
\end{Theorem}

\begingroup
\postdisplaypenalty=10000
\renewcommand{\endProof}{\par\nobreak\hfill$\blacksquare$\par\vspace{0.1cm}}
\begin{Proof}
Let $\{\phi_i(\bm{x})\}$ be an orthonormal basis for $L^2(\Omega)$. Then, using the methods of Section~\ref{subsec:distribution-direct-product}, one can show that $\{\phi_i(\bm{x})\phi_j(\bm{y})\}$ is a basis for $L^2(\Omega\times\Omega)$. Expanding $k$ with respect to this basis gives us
\begin{equation}\label{eq:hilbert-schmidt-kernel-product-basis-expansion}
k(\bm{x},\bm{y})=\sum_{i,j=1}^\infty k_{ij}\phi_i(\bm{x})\phi_j(\bm{y})
\end{equation}
where the convergence of the sum is in the $L^2(\Omega\times\Omega)$ norm and
\begin{equation}\label{eq:hilbert-schmidt-kernel-product-basis-coefficients}
k_{ij}:=\int_\Omega\int_\Omega k(\bm{x},\bm{y})\phi_i(\bm{x})\phi_j(\bm{y})\,\mathrm{d}\bm{x}\,\mathrm{d}\bm{y}.
\end{equation}
Furthermore, by \eqref{eq:hilbert-basis-parseval-characterization} we have
\begin{equation}\label{eq:hilbert-schmidt-kernel-coefficient-square-sum}
\int_\Omega\int_\Omega|k(\bm{x},\bm{y})|^2\,\mathrm{d}\bm{x}\,\mathrm{d}\bm{y}=\sum_{i,j=1}^\infty|k_{ij}|^2.
\end{equation}
We now define the operator $K_n\in\mathcal{L}(L^2(\Omega))$ by
\begin{equation}\label{eq:hilbert-schmidt-truncated-kernel-operator}
K_n(u)(\bm{x}):=\int_\Omega k_n(\bm{x},\bm{y})u(\bm{y})\,\mathrm{d}\bm{y}
\end{equation}
where
\begin{equation}\label{eq:hilbert-schmidt-rectangular-kernel-truncation}
k_n(\bm{x},\bm{y})=\sum_{i,j=1}^n k_{ij}\phi_i(\bm{x})\phi_j(\bm{y}).
\end{equation}
We refer to $k_n$ and $K_n$ as \textbf{separable kernels} and \textbf{separable operators}, respectively. It is easy to see that a separable operator has finite rank and is thus compact.

We now use the techniques of Lemma~\ref{lem:hilbert-schmidt-kernel-boundedness} to get
\begin{equation}\label{eq:hilbert-schmidt-truncation-operator-norm-bound}
\|K-K_n\|^2\leq\int_\Omega\int_\Omega|k(\bm{x},\bm{y})-k_n(\bm{x},\bm{y})|^2\,\mathrm{d}\bm{x}\,\mathrm{d}\bm{y}.
\end{equation}
Now we use \eqref{eq:hilbert-schmidt-kernel-coefficient-square-sum} to get
\begin{equation}\label{eq:hilbert-schmidt-truncation-coefficient-tail}
\begin{aligned}
&\lim_{n\to\infty}\int_\Omega\int_\Omega|k(\bm{x},\bm{y})-k_n(\bm{x},\bm{y})|^2\,\mathrm{d}\bm{x}\,\mathrm{d}\bm{y}\\
&\qquad=\lim_{n\to\infty}\sum_{i,j=n+1}^\infty|k_{ij}|^2=0.
\end{aligned}
\end{equation}
Thus, $K_n$ converges to $K$ in the operator norm, so Theorem~\ref{thm:norm-limit-of-compact-operators} implies that $K$ is compact.
\end{Proof}
\endgroup
